\documentclass{article}
\usepackage[T1]{fontenc}
\usepackage{lmodern}
\usepackage{graphicx}
\usepackage{amsmath,amssymb}
\usepackage[a4paper,margin=2cm]{geometry}
\usepackage[absolute,overlay]{textpos}
\setlength{\TPHorizModule}{1mm}
\setlength{\TPVertModule}{1mm}
\usepackage[indonesian]{babel}
\usepackage{hyperref}
\setlength{\emergencystretch}{2em}
% unit: Halaman 26

\begin{document}

% === Halaman 26 (python/native) ===

26

• Bit-bit aliran kunci untuk enkripsi/dekripsi disebut keystream

• Keystream dibangkitkan oleh keystream generator berdasarkan umpan (seed) U.

• Enkripsi dan dekripsi dengan cipher alir komputasinya sederhana dan cepat 

• Keystream c1, c2, … di-XOR-kan dengan aliran bit-bit plainteks, p1, p2, …, menghasilkan aliran bit-bit 

cipherteks c1, c2, …:



ci = pi  ki 

• Di sisi penerima dibangkitkan keystream yang sama untuk mendekripsi aliran bit-bit cipherteks c1, 

c2, … menjadi plainteks, p1, p2, …:



pi = ci  ki 

Keystream

Generator

Keystream

Generator





ki

ki

pi

pi

ci

Enkripsi

Dekripsi

Keystream

Keystream

Plainteks

Plainteks

Cipherteks

Pengirim

Penerima

Diagram Umum Cipher alir

\end{document}
